\par\Needspace{4\baselineskip}
\begin{center}
\hypertarget{native-033-article015}{}\pdfbookmark[2]{15.}{article-015}15.
\end{center}
Da die mittlere Anomalie $M$, wie wir gesehen haben, durch $v$ und $\varphi$ vollständig bestimmt ist, ebenso wie $v$ durch $M$ und $\varphi$, müssen die differenziellen Änderungen dieser drei Größen offenbar eine Bedingungsgleichung erfüllen, wenn man alle drei zugleich als veränderlich betrachtet. Es lohnt sich, diese Gleichung herzuleiten. Differenzieren wir \SourceLeaf{034}zunächst Gleichung VII aus Art. 8, so erhalten wir $\frac{\mathrm{d}E}{\sin E}=\frac{\mathrm{d}v}{\sin v}-\frac{\mathrm{d}\varphi}{\cos\varphi}$. Differenzieren wir ebenso Gleichung XII, so ergibt sich $\mathrm{d}M=(1-e\cos E)\mathrm{d}E-\sin E\cos\varphi\,\mathrm{d}\varphi$. Eliminieren wir $\mathrm{d}E$ aus diesen Differenzialgleichungen, so erhalten wir
\[
\mathrm{d}M=\frac{\sin E(1-e\cos E)}{\sin v}\,\mathrm{d}v-\left(\sin E\cos\varphi+\frac{\sin E(1-e\cos E)}{\cos\varphi}\right)\mathrm{d}\varphi
\]
oder, wenn wir für $\sin E$ und $1-e\cos E$ die jeweiligen Werte aus den Gleichungen VIII und III einsetzen,
\[
\mathrm{d}M=\frac{r r}{a a\cos\varphi}\,\mathrm{d}v-\frac{r(r+p)\sin v}{a a\cos\varphi^2}\,\mathrm{d}\varphi
\]
oder schließlich, wenn wir beide Koeffizienten allein durch $v$ und $\varphi$ ausdrücken,
\[
\mathrm{d}M=\frac{\cos\varphi^3}{(1+e\cos v)^2}\,\mathrm{d}v-\frac{(2+e\cos v)\sin v\cos\varphi^2}{(1+e\cos v)^2}\,\mathrm{d}\varphi.
\]
Betrachten wir umgekehrt $v$ als Funktion der Größen $M$ und $\varphi$, so nimmt die Gleichung folgende Form an
\[
\mathrm{d}v=\frac{a a\cos\varphi}{r r}\,\mathrm{d}M+\frac{(2+e\cos v)\sin v}{\cos\varphi}\,\mathrm{d}\varphi
\]
oder, wenn wir $E$ anstelle von $v$ einführen\,\raisebox{-.35ex}{.}
\[
\mathrm{d}v=\frac{a a\cos\varphi}{r r}\,\mathrm{d}M+\frac{a a}{r r}(2-e\cos E-e e)\sin E\,\mathrm{d}\varphi.
\]

\par\Needspace{4\baselineskip}
\begin{center}
\hypertarget{native-034-article016}{}\pdfbookmark[2]{16.}{article-016}16.
\end{center}
Für die vollständige Bestimmung des Radiusvektors $r$ aus $v$ und $\varphi$ oder aus $M$ und $\varphi$ benötigt man zusätzlich $p$ oder $a$. Sein Differenzial besteht also aus drei Gliedern. Durch Differenzieren von Gleichung II aus Art. 8 erhalten wir
\[
\frac{\mathrm{d}r}{r}=\frac{\mathrm{d}p}{p}+\frac{e\sin v}{1+e\cos v}\,\mathrm{d}v-\frac{\cos\varphi\cos v}{1+e\cos v}\,\mathrm{d}\varphi.
\]
Setzen wir hier $\frac{\mathrm{d}p}{p}=\frac{\mathrm{d}a}{a}-2\operatorname{tang}\varphi\,\mathrm{d}\varphi$ (was aus der Differenzierung von Gleichung I folgt) und drücken wir gemäß dem vorigen Artikel $\mathrm{d}v$ durch $\mathrm{d}M$ und $\mathrm{d}\varphi$ aus, so erhalten wir nach den nötigen Vereinfachungen
\[
\frac{\mathrm{d}r}{r}=\frac{\mathrm{d}a}{a}+\frac{a}{r}\operatorname{tang}\varphi\sin v\,\mathrm{d}M-\frac{a}{r}\cos\varphi\cos v\,\mathrm{d}\varphi
\]
oder
\[
\mathrm{d}r=\frac{r}{a}\,\mathrm{d}a+a\operatorname{tang}\varphi\sin v\,\mathrm{d}M-a\cos\varphi\cos v\,\mathrm{d}\varphi.
\]

Diesen Formeln liegt ebenso wie den im vorigen Artikel hergeleiteten die Annahme zugrunde, dass $v$, $\varphi$ und $M$ oder vielmehr $\mathrm{d}v$, $\mathrm{d}\varphi$ und $\mathrm{d}M$ in Teilen des Radius \SourceLeaf{035}ausgedrückt werden. Will man also die Änderungen der Winkel $v$, $\varphi$ und $M$ in Bogensekunden ausdrücken, so muss man entweder diejenigen Teile der Formeln, in denen $\mathrm{d}v$, $\mathrm{d}\varphi$ oder $\mathrm{d}M$ vorkommen, durch $206264{,}8$ teilen oder diejenigen, die $\mathrm{d}r$, $\mathrm{d}p$ oder $\mathrm{d}a$ enthalten, mit derselben Zahl multiplizieren. Die Formeln des vorigen Artikels, die in dieser Hinsicht homogen sind, bedürfen also keiner Änderung.
